
\section{Discusión}

Los resultados sugieren una reformulación general de la reconstrucción constructiva computacional.

El objeto terminal puede preservar información en múltiples canales:

\begin{equation}
\boxed{
\text{geometría}
+
\text{microestructura}
+
\text{contactos}
+
\text{carga}
+
\text{vacío}
+
\text{procedencia}
+
\text{daño}
+
\text{sitio}.
}
\end{equation}

Sin embargo, el estudio adversarial demuestra que esa información no debe sobreinterpretarse. El terminal no contiene toda la historia. Contiene restricciones sobre la historia.

La tesis defendible es:

\begin{quote}
un estado terminal multimodal, combinado con leyes físicas y evidencia contextual explícita, puede reducir el espacio de historias constructivas y, en determinados casos, identificar clases funcionales de procesos temporales que no sobreviven en el objeto final.
\end{quote}

Los cinco controles también muestran que lo que parece un problema de fuerza puede ser un problema de representación:

\begin{itemize}
\item Golden Gate: accesibilidad y secuenciación;
\item Eiffel: modularidad, ajuste y fijación temporal;
\item Hoover: hidráulica, térmica y estado del material;
\item Empire State: pipeline y logística;
\item Sydney: geometría generativa y soporte temporal.
\end{itemize}

Esto justifica el principio de no introducir nueva física antes de agotar ontología y procesos convencionales.

\section{Conclusiones}

Se ha construido un marco formal en el que:

\begin{enumerate}
\item el presente se separa del estado terminal as-built;
\item la cronología se representa como orden parcial y no como altura;
\item materia y espacio negativo se modelan conjuntamente;
\item los vacíos poseen genealogía y escala topológica;
\item mecánica y acceso actúan como gates;
\item procedencia y logística restringen historias;
\item Bayes preserva equifinalidad e incertidumbre;
\item el motor contrafactual convierte historias en predicciones;
\item los objetos temporales desaparecidos pueden inferirse a nivel funcional;
\item la identificabilidad se comprueba antes de afirmar un mecanismo;
\item una fuerza residual se calcula sólo después de agotar explicaciones convencionales.
\end{enumerate}

Los cinco benchmarks no validan todavía el método como sistema automático, pero muestran consistencia de representación en dominios constructivos muy diferentes. El estudio adversarial es igualmente importante: demuestra que geometría aislada es insuficiente y obliga a degradar afirmaciones inicialmente demasiado fuertes.

El objetivo científico final es

\begin{equation}
\boxed{
P(
\text{historias constructivas}
\mid
\text{objeto terminal},
\text{evidencia},
\text{leyes},
\text{contexto}
).
}
\end{equation}

El valor de IRL no dependerá de producir una explicación extraordinaria de Khufu. Dependerá de eliminar historias, declarar no identificabilidad, anticipar evidencia nueva y sobrevivir controles prospectivos donde la respuesta permanezca realmente oculta.
